% GERADO por scripts/analyse_feedback.py — NÃO EDITAR À MÃO.
% Correr o script outra vez depois de a recolha fechar; o texto acompanha os dados.
\subsection{Perceived usefulness of the alerts delivered}
\label{sec:av_feedback}

The channel now accompanies every alert delivered with two buttons, \emph{useful} and \emph{did not help}, and records the vote of whoever presses. The analysis rules were fixed before a single vote existed and were not changed afterwards: a minimum of twenty effective votes to report any proportion, one vote per person and per alert with the last replacing the previous one, a safeguard that repeats the computation without the dominant voter when that voter accounts for more than forty per cent of the votes, always subject to the same minimum, and Wilson confidence intervals, calculated without adjustment for clustering.
Only votes on alerts present in the shared log were counted.

131 valid votes were recorded between 2026-09-02 and 2026-09-11, corresponding to 91 effective votes from 3 distinct people, on 62 alerts. 11 votes were later changed by the same person, and only the last of each pair counts. Another 29 repeat an earlier vote without changing it and count only once. A further 5 votes, not counted among the 131 above, were excluded for not corresponding to any alert in the shared log.

Of the 91 effective votes, 86 rated the alert as useful, that is 95\%, with a 95\% Wilson confidence interval between 88\% and 98\%. This binomial calculation does not adjust for dependence between votes by the same person or on the same alert; its width does not represent all the uncertainty in this sample.

A single person accounts for 58\% of the effective votes, above the limit of 40\% fixed in the protocol. Excluding that person, 38 votes remain, of which 34 rate the alert as useful.
The proportion without that person is 89\%, with an interval between 76\% and 96\%. This is the reading to retain.

The independence between whoever builds the system and whoever rates it is a condition of this sample: the author is not among the votes counted. The votes come from readers of the public channel, and the log keeps of each only the conversation identifier assigned by the platform.

These votes are observational feedback in a real context and do not replace the controlled study that Section~\ref{sec:con_limitacoes} identifies as the gap of greatest weight in this work, and which remains to be carried out.

Four limitations accompany this result and none of them is resolved by more votes. The first is self-selection: whoever wants to votes, and those who find an alert indifferent tend to press nothing, which pushes the sample towards the extremes. The second is the absence of a counterfactual, since no group received the price change without the explanation that accompanies it; nothing here attributes the usefulness to the explanation itself. The third is the distance between perceived usefulness and a better decision, which is the founding hypothesis of this work and remains untested. The fourth is not knowing who votes: the channel is public and it is not known whether those who answer belong to the audience the dissertation assumes.
